\section{Stability of the horizontal weights}\label{c7:stability}

The horizontal representation has signed and complex weights. Its use requires an absolute estimate: configurations with many particles or a large common charge must have a summable total weight. We prove that estimate directly from the signed height lists. The estimate applies both to the absolute Gaussian moment before collisions and to the Wick-evaluated gas after collisions.

For the post-collision gas, let \(W_{d,\theta}\) be the weight \eqref{c7:horizontal-weight}, with the deterministic \(B_T\)-powers omitted. Thus it includes the local constants, the state factor, and every forward residual pair factor, including nonzero self-images. We next define the precise real product needed before collision.

For a basic label \(v\) at a fixed \(-1/2<d<0\), let \(\widetilde{\mathscr H}_v\) be obtained from its signed \(H\)-list by replacing every entry \((a,r)\) by \((a,r/2),(-a,r/2)\), retaining multiplicities. Evaluate the basic short and outer lists at this same \(d\). Define \(\gamma_v^{\rm pre}\) to be the modulus of the right side of \eqref{c7:local-factor-rule} with its internal \(H\)-list replaced by \(\widetilde{\mathscr H}_v\), leaving those short and outer lists unchanged. Similarly, for \(x>0\), define \(\mathscr R^{\rm pre}_{uv}(x)\) as the modulus of the right side of \eqref{c7:forward-residual} with both \(H\)-lists replaced by their tilded lists, again leaving the short and outer lists at the same \(d\). Thus the internal \(H\) self-list is symmetrized as well as the external pairs. For example, for \(U\) with \(a=3/2+d\), its internal \(H\)-logarithm is
\[
\frac14\bigl(K_{\rm diag}(0)+K_{\rm diag}(2ia)\bigr),
\]
which is finite because \(|2a|<3\); it is generally different from the signed local logarithm \(K_{\rm diag}(0)/2\).

For a configuration \(\mathcal Y=((x_j,v_j))_{j=1}^N\) of basic labels and a compatible state path \(p\), define, for real \(\vartheta\),
\begin{equation}\label{j:pre-weight}
\begin{aligned}
W^{\rm pre}_{d,\vartheta}(\mathcal Y,p)
:={}&\exp\!\left[-\pi\int_T
(p+(\vartheta/(2\pi))\mathbf1)^t\mathcal L
(p+(\vartheta/(2\pi))\mathbf1)\,dx\right]\\
&\times\prod_j\gamma_{v_j}^{\rm pre}
\prod_{\substack{1\le j,k\le N,\ n\in\mathbb Z\\x_k+nT>x_j}}
\mathscr R^{\rm pre}_{v_jv_k}(x_k+nT-x_j).
\end{aligned}
\end{equation}
This is a nonnegative product and has no Berry factor. The common scalar prefactors from the real heat completion are outside this configuration weight.

\begin{theorem}[Absolute stability at a fixed displacement]\label{j:absolute-stability}
There is \(0<d_0<1/2\) with the following property. Fix a compact set \(\Theta\subset\mathbb C\), and fix either a pre-collision displacement \(-d_0<d<0\) or a post-collision displacement \(0<d<d_0\). There are constants \(c>0\) and \(C<\infty\), depending on \(d\) and \(\Theta\), such that the following holds for every positive integer \(T\) and \(\theta\in\Theta\).

Consider a finite configuration \(\mathcal Y\) with distinct real positions on \(\mathbb R/T\mathbb Z\), using the six basic labels in the pre-collision case and those labels together with \(F_{\rm pair}\) in the post-collision case. Give it a state path \(p\) with jumps \(-q_i\) and closure in \(\mathbb Z^2/8\mathbb Z(1,-1)\). Let \(n_j\) be the number of physical particles in \([j,j+1)\), and let
\[
p_{c,j}=\frac{p_1(j+)+p_2(j+)}{\sqrt2},
\qquad 0\le j<T.
\]
Then
\[
\left|W_{d,\theta}(\mathcal Y,p)\right|
\ \le\
\exp\!\left(CT-c\sum_{j=0}^{T-1}(n_j^2+p_{c,j}^2)\right)
\]
in the post-collision case. The same bound holds for \(W^{\rm pre}_{d,\Re\theta}(\mathcal Y,p)\) in the pre-collision case.
\end{theorem}

For the actual pre-collision Gaussian majorant, write \(\theta=\vartheta+i\tau\) and
\[
Z=\sum_j\sum_{(a,r)\in\mathscr H_{v_j}}ra\,e_{\sigma(v_j)}
=\Im X,
\]
using the original \(H\)-lists. The explicit factor from \eqref{c7:absolute-heat} that lies outside \eqref{j:pre-weight} is
\[
\exp\!\left(\frac{T\tau^2|\mathbf1|^2}{4\pi}
+\tau\mathbf1\cdot Z\right)
\le \exp(C_\Theta(T+N)).
\]
The inequality follows from bounded heights and charges per basic particle. This factor is absorbed by the theorem's quadratic count bound. The \(B_T\)-powers are outside both weights as well. They have modulus at most one on the post-collision contours: their simplified one-particle factors use \(B(v+i)B(v-i)=B(v)\), while \(|B(v)|<1\) for \(|\Im v|<3/2\) and \(|B(v)|>1\) just beyond either boundary line; the fused factor is one. Before collision they have a finite bound per particle at fixed \(T,m,d\), which gives a further absorbable linear count cost. These observations give the Gaussian integration at each fixed period and finite site count. The later marked application will use the separate proof at \(d=1/8\) in Section~\ref{c7:shifted}.

\subsection{The matrix of signed height measures}

Write \(\mu_i^s,\mu_i^o,\mu_i^H\) for the signed measures on real heights in the short, outer, and \(H\) columns of \eqref{j:particle-list}, for a label \(i\) in one species. For example,
\[
\mu_P^s=\delta_{1/2},\qquad
\mu_P^o=\delta_1,\qquad
\mu_P^H=\delta_0.
\]
The separate coefficient sign in the table belongs to the local factor and does not multiply these measures. The fused measures are the union of the \(U,L\) measures at zero displacement:
\[
\mu_{F_{\rm pair}}^s=-2\delta_0,\qquad
\mu_{F_{\rm pair}}^o=\delta_{1/2}-\delta_{-1/2},\qquad
\mu_{F_{\rm pair}}^H=-\delta_{3/2}-\delta_{-3/2}.
\]
In the order \(P,M,O_+,O_-,U,L,F_{\rm pair}\), the total of both the short measure and the \(H\)-measure is
\[
\bar q_i=(1,-1,0,0,-1,-1,-2)_i
\]
in that species. Thus the physical charge vector is \(q_{(\alpha,i)}=\bar q_i e_\alpha\) for species \(\alpha\). For a label \(\iota=(\alpha,i)\), put
\[
c_\iota=\frac{\bar q_i}{\sqrt2},\qquad
\boldsymbol c=(c_\iota)_\iota,\qquad a_c=\frac89.
\]
The jump of the common charge at label \(\iota\) is \(-c_\iota\).

For \(k>0\), the short and outer height kernels are
\[
\begin{aligned}
p(a)&=\frac{\cosh((1/2-|a|)k)}{2\sinh(k/2)}
       &&(|a|\le1,\ \text{period }1),\\
o(a)&=\frac{\sinh((1-|a|)k)}{2\cosh k}
       &&(|a|\le2,\ \text{antiperiod }2).
\end{aligned}
\]
They are positive covariance kernels. More explicitly, their Green series are
\[
p(a)=\sum_{n\in\mathbb Z}\frac{k e^{2\pi ina}}{k^2+(2\pi n)^2},
\qquad
o(a)=\frac12\sum_{n\in\mathbb Z}
\frac{k e^{\pi i(n+1/2)a}}{k^2+\pi^2(n+1/2)^2}.
\]
Set
\[
C_i(k)=\int\cosh(ka)\,d\mu_i^H(a),\qquad
S_i(k)=\int\sinh(ka)\,d\mu_i^H(a).
\]
Define the dimensionless matrix \(\mathcal H\) by
\begin{equation}\label{j:height-matrix}
\begin{aligned}
\mathcal H_{\alpha i,\beta j}(k)
={}&\delta_{\alpha\beta}
\left[
\iint p(a-b)\,d\mu_i^s(a)d\mu_j^s(b)
+\iint o(a-b)\,d\mu_i^o(a)d\mu_j^o(b)
\right]\\
&+\sum_{\sigma=\pm}(P_\sigma)_{\alpha\beta}\Phi^\sigma_{ij}(k),
\end{aligned}
\end{equation}
where the two uses have
\[
\begin{aligned}
\Phi^{\sigma,{\rm pre}}_{ij}
&=-\mathsf h_\sigma C_iC_j,\\
\Phi^{\sigma,{\rm post}}_{ij}
&=-\mathsf h_\sigma(C_iC_j-S_iS_j)
+4\mathbf1_{\{\sigma=-\}}\iint
\sinh\!\big(k(|a-b|-3)_+\big)\,d\mu_i^H(a)d\mu_j^H(b).
\end{aligned}
\]
The first formula is the Gaussian absolute moment: reality of \(\xi\) on the real line gives
\(\Re\xi(x+ia)=(\xi(x+ia)+\xi(x-ia))/2\). The second uses the original charges and
\(\cosh(k(a-b))=\cosh(ka)\cosh(kb)-\sinh(ka)\sinh(kb)\).

We verify the continuation term, since it controls the collision directions. The only singular part of \(K_0\) below height \(4\) is
\(K_{\rm lead}(s)=-4P_-\log(1+s^2/9)\). For \(k>0\), the transform of its negative real part at height \(a\) is
\[
\frac{k}{2\pi}\widehat{-\Re K_{\rm lead}(x+ia)}
=P_-\left[-4e^{-3k}\cosh(ak)
+4\sinh\!\big((|a|-3)_+k\big)\right].
\]
Indeed the two real logarithms involve \(x^2+(3-a)^2\) and \(x^2+(3+a)^2\); their transforms use \(e^{-|3-a|k}\) and \(e^{-|3+a|k}\). Crossing height \(3\) therefore replaces the growing factor \(e^{-(3-|a|)k}\) by \(e^{-|3-|a||k}\), giving the displayed correction. In the pre-collision absolute moment all extreme height differences are less than \(3\), so that correction is absent.

Let \(V^r_{\iota\jmath}(x)\) be the even extension of the negative log magnitude of one forward residual pair, for \(x\ne0\). In the pre-collision use it denotes the analogous real symmetrized pair. It is integrable: possible singularities at zero are logarithmic, and the residual factors decay exponentially at infinity. Expanding the short logarithm gives transform \((2\pi/k)(p(a)-1/k)\); the outer logarithm gives \((2\pi/k)o(a)\). The relation \(K^r=K_0+\pi|x|\mathcal S\) adds \((2\pi/k)\mathcal S/k\) to the \(H\) residual transform. Since both charge totals are \(\bar q_i\), removing the full-factor linear terms leaves
\begin{equation}\label{j:residual-symbol}
\widehat V^r(k)=\frac{2\pi}{k}
\left(\mathcal H(k)-\frac{a_c}{k}\boldsymbol c\boldsymbol c^t\right),
\qquad k>0.
\end{equation}
There is no constant ambiguity here because the residual kernel decays at infinity.

\subsection{Coercivity at a fixed displacement}

The useful positive part of \(\mathcal H\) is most easily found after normalizing the fused coordinate. In this calculation only, put \(F=F_{\rm pair}/2\), halving all three signed measures. It can be included as an auxiliary coordinate in the pre-collision calculation and removed afterward. Let
\[
D_0=\operatorname{diag}(1,1,1,1,1,1,2)
\]
on \(P,M,O_+,O_-,U,L,F\), and let \(D_F=\operatorname{diag}(D_0,D_0)\) on the two species. Then
\(\mathcal H_{\rm phys}=D_F\mathcal H_{\rm norm}D_F\). A coefficient vector \(v_{\rm phys}\) has normalized coordinates \(D_Fv_{\rm phys}\), so physical particle count is represented there by
\[
\mathfrak n_{\rm norm}\big|_{\text{one species}}
=D_0^{-1}\mathbf1_7=(1,1,1,1,1,1,\tfrac12).
\]
The normalized charge row is \((1,-1,0,0,-1,-1,-1)\). In particular \(F\) has count \(1/2\) and charge \(-1\), whereas \(R_d=(U+L)/2\) has count \(1\) and charge \(-1\).

This distinction explains the fixed-displacement hypothesis. At \(d=0\), the height measures of \(R_d\) and \(F\) agree but their count values do not. Thus the all-ones count bound below cannot be uniform through \(d=0\). The same-species \(U,L\) factor also has the simple collision pole from Section~\ref{c7:boson}, with magnitude proportional to \(1/|x\pm2id|\) near collision. The pointwise estimate consequently needs constants depending on \(d\). The contour argument sums the series at fixed endpoints on the two sides of the collision.

In normalized coordinates subtract from \(\mathcal H\) the species-off-diagonal matrix
\[
K_{\rm norm}(k)=
\begin{pmatrix}0&\lambda VV^t\\ \lambda VV^t&0\end{pmatrix},
\qquad
\lambda=\tfrac38\tanh(2k),
\]
where
\[
V_P=0,\quad V_M=1,\quad V_{O_\pm}=\tfrac12,\quad
V_U=V_L=e^{-|d|k},\quad V_F=1.
\]
Write \(\mathcal H'=\mathcal H-K\), using congruence by \(D_F\) for both matrices when returning to the physical basis. Reflection of heights, together with negation of outer signs, preserves these forms; the outer at height \(1\) is unchanged by antiperiodicity. Thus \(P,M,F\) are even, while \(O_+\) and \(O_-\), and \(U\) and \(L\), are exchanged. The exact measure identity
\[
O_++O_-=P+M
\]
also holds for \(V\). Its two species copies generate a fixed nullspace \(\mathcal N\). Both physical count and common charge vanish on \(\mathcal N\).

\begin{lemma}[Coercive matrix part]\label{j:stability-coercivity}
For each fixed displacement in Theorem~\ref{j:absolute-stability}, the physical matrix \(\mathcal H'(k)\) has nullspace \(\mathcal N\) for every \(k>0\). On the quotient by \(\mathcal N\), it is bounded below near zero by a positive constant times
\[
kI+\frac{\boldsymbol c\boldsymbol c^t}{k},
\]
and it converges exponentially at infinity to a strictly positive matrix \(\mathcal H'(\infty)\). The physical limit \(Q_\infty=\mathcal H(\infty)\) is entrywise nonnegative.

Consequently one can choose \(0<e<2\) and \(\gamma>0\), depending on the fixed displacement, such that
\begin{equation}\label{j:stability-J}
J(k):=\mathcal H'(k)-\mathcal H'(\infty)\tanh(ek),
\qquad
\frac{J(k)}{k}\ge
\gamma e^{-k^2}\left(
\mathfrak n\mathfrak n^t+
\frac{\boldsymbol c\boldsymbol c^t}{k^2}\right),\quad k>0,
\end{equation}
where \(\mathfrak n\) is the all-ones vector on the actual physical labels. In the pre-collision case the matrices in this statement are restricted to the six basic labels.
\end{lemma}

\begin{proof}
We first analyze the even directions, then the odd directions, in each species channel \(\sigma=\pm\).

\paragraph{The even core.}
In the quotient, the even direction \((O_++O_-)/2\) equals \((P+M)/2\). On \(P,M,F\), the short and outer measures give
\begin{equation}\label{j:stability-core}
\begin{gathered}
A_\sigma=M_0-\mathsf h_\sigma XX^t-\sigma\lambda VV^t,\\
M_0=
\begin{pmatrix}
f+g&-f+g&-l+j_0\\
-f+g&f+g&l+j_0\\
-l+j_0&l+j_0&f+g/2
\end{pmatrix},
\end{gathered}
\end{equation}
with
\[
\begin{gathered}
f=p(0),\quad g=o(0),\quad l=p(1/2),\quad j_0=o(1/2),\\
X=(1,1-2\cosh k,-\cosh(3k/2))^t,\quad V=(0,1,1)^t.
\end{gathered}
\]
For example, the normalized \(F\) short is \(-\delta_0\) and its two outers have weights \(1/2,-1/2\), producing the third row of \(M_0\). The three even \(H\)-moments give \(X\). The same core submatrix is present for nonzero \(d\). Strict Green positivity makes \(M_0>0\): cancellation of its short measure requires equal \(P,M\) coefficients and zero \(F\) coefficient, and the remaining outer measure then forces both coefficients to vanish.

Let \(u=e^{-k/2}\) and set
\[
x=X^tM_0^{-1}X,\qquad y=X^tM_0^{-1}V,\qquad v=V^tM_0^{-1}V.
\]
Direct inversion gives
\[
\begin{aligned}
x&=\frac{(1-u^2)(u^4+1)(u^8+u^6-u^4+u^2+1)}
        {3u^6(u^2+1)},\\
y&=\frac{2(u^2-1)(u^4+1)}{3u^3},\qquad
v=\frac{7u^2-12u+7}{3(1-u^2)}.
\end{aligned}
\]
The signs of \(1-\mathsf h_+x\) and \(1-\mathsf h_-x\) are respectively positive and negative. After positive denominators are removed their numerators are
\[
(u^4+1)(3u^8+2u^6+2u^4+2u^2+3),\quad
-(1-u^2)^2(u^4+1)(u^8+2u^6-2u^4+2u^2+1).
\]
The last parenthesis is positive: with \(v_0=u^2\) it equals
\((1-v_0^2)^2+2v_0(1+v_0^2)\). Put
\[
w_\sigma=v+\frac{\mathsf h_\sigma y^2}{1-\mathsf h_\sigma x}.
\]
Substitution of the rational formulas for \(\mathsf h_\sigma(t=u^2)\) and
\(\lambda=3(1-u^8)/(8(1+u^8))\) gives
\[
\frac1{w_\sigma}-\sigma\lambda
=\frac{1-u^2}{u}\frac{N_\sigma(z)}{D_\sigma(z)}>0,
\qquad z=u+u^{-1}-2.
\]
The coefficients, in ascending order, are
\[
\begin{aligned}
N_+&:(96,1280,5496,9312,7696,3360,772,84,3),\\
D_+&:(128,2240,13888,38304,55392,45584,21984,6128,912,56),\\
N_-&:(32,256,424,1024,1384,832,232,28,1),\\
D_-&:(128,2240,13760,36704,49120,36176,15264,3664,464,24).
\end{aligned}
\]
Clearing denominators verifies the identities, and all these coefficients are positive for \(z\ge0\). The finite positive right side prevents \(w_\sigma\) from vanishing. In the plus channel, \(M_0-\mathsf h_+XX^t>0\), so \(w_+>0\), and the displayed inequality is \(1-\lambda w_+>0\). Both negative rank-one updates preserve positivity. In the minus channel, the first update has exactly one negative direction. Since \(w_-\to-3\) as \(u\to0\), it is negative throughout; the inequality now gives \(1+\lambda w_-<0\). The determinant lemma makes the determinant positive after addition of \(\lambda VV^t\). A positive rank-one update cannot add negative directions, so this removes the unique negative direction. Hence \(A_\sigma>0\) for every \(k>0\).

The endpoint expansions give uniform core gaps. With \(q=(1,-1,-1)^t\),
\[
M_0=\frac{qq^t}{k}
+\frac{k}{24}\begin{pmatrix}14&10&7\\10&14&5\\7&5&8\end{pmatrix}
+O(k^3),\qquad
X=q+k^2(0,-1,-9/8)^t+O(k^4).
\]
They yield
\[
\frac{A_-}{k}\longrightarrow
\begin{pmatrix}2&0&0\\0&3/4&1/4\\0&1/4&1/4\end{pmatrix},
\]
while \(A_+\) has leading term \((8/9)qq^t/k\). Its order-\(k\) restriction to \(q^\perp\), in the basis \((1,1,0),(1,0,1)\), is
\(\left(\begin{smallmatrix}5/4&3/4\\3/4&3/4\end{smallmatrix}\right)>0\).
At infinity,
\[
A_\sigma^\infty
=\operatorname{diag}(1,1,3/4-\mathbf1_{\{\sigma=-\}})
-\sigma(3/8)VV^t.
\]
The leading principal minors are \((1,5/8,3/32)\) for \(+\) and \((1,11/8,1/32)\) for \(-\). Positivity on the intervening compact frequency ranges therefore gives
\[
A_\sigma\ge c_*\min(k,1)I,
\]
with an additional \(c_*qq^t/k\) term for \(A_+\) near zero.

\paragraph{The displaced directions.}
The remaining even direction is \(R_d=(U+L)/2\); put \(Z_d=R_d-F\) and \(D_*=|d|k\). The \(H\)-moment of \(R_d\) is \(-\cosh((3/2+d)k)\). The odd directions are
\[
O_{\rm odd}=(O_+-O_-)/2,\qquad W_d=(U-L)/2,
\]
whose \(H\)-sinh moments are \(-\sinh k\) and \(-\sinh((3/2+d)k)\). Write \(\mathcal T\) for the continued correction in the post-collision minus channel. Summing the opposite extreme charges with their actual half weights gives precisely
\begin{equation}\label{j:height-continuations}
\mathcal T_{R_dR_d}=2\sinh(2D_*),\qquad
\mathcal T_{R_dF}=2\sinh D_*,\qquad
\mathcal T_{W_dW_d}=-2\sinh(2D_*),
\end{equation}
and no other continued entries. These terms are absent before collision.

When \(D_*\) and \(|d|\) are sufficiently small, the square of \(Z_d\) is
\(D_*(1+O(D_*+|d|))\) and its crosses with the core are \(O(D_*)\), uniformly. For the field part, its short square is exactly
\[
\tfrac32p(0)+\tfrac12p(2d)-2p(d)
=f(\cosh D_*-1)^2+\sinh D_*-\tfrac14\sinh(2D_*),
\]
and its outer square is
\[
o(0)-o(d)+o(1+d)-\tfrac12o(1+2d).
\]
Each has leading term \(D_*/2\); the formulas for \(p,o\) also give the asserted cross bounds. At small \(k\), the potentially larger second-order bounds \(O(D_*^2/k)\) are still \(O(|d|D_*)\). The uncontinued \(H\)-square difference is \(O(D_*^2)\) and its crosses are \(O(D_*)\), using
\[
\mathsf h_+=\frac1{9k}+O(k),\quad
\mathsf h_-=\frac1k+O(k)\quad(k\downarrow0),\qquad
\mathsf h_+=O(e^{-4k}),\quad
\mathsf h_-=4e^{-3k}+O(e^{-4k})\quad(k\to\infty).
\]
The continued contribution to the \(Z_d\) square is
\(2\sinh(2D_*)-4\sinh D_*=O(D_*^3)\). The \(V\)-terms have the same bounds. The Schur cost of all crosses is at most
\[
\frac{C D_*^2}{\min(k,1)}
\le C\max(|d|,D_*)D_*
\]
when \(D_*\) and \(|d|\) are small, because
\(D_*/\min(k,1)\le\max(|d|,D_*)\). Absorption against the core gap leaves a positive multiple of the \(D_*\) square and part of the plus charge gap. The direction \(Z_d\) has zero total charge.

For large \(k\) with \(D_*\) bounded below by a fixed positive constant, the core cross row of \(R_d\) is
\[
e^{-D_*}A^\infty_{\sigma,F\bullet}+O(e^{-c'k}),
\]
and its square is
\[
e^{-2D_*}A^\infty_{\sigma,FF}
+\tfrac12(1-e^{-2D_*})+O(e^{-c'k}).
\]
The field square tends to \(1/2+e^{-2D_*}/4\), and the field \(R_d,F\) entry to \(3e^{-D_*}/4\). In the minus channel the leading uncontinued height terms are \(-e^{2dk}\) and \(-e^{dk}\). For \(d>0\), \eqref{j:height-continuations} changes these exactly to
\[
-e^{2D_*}+2\sinh(2D_*)=-e^{-2D_*},\qquad
-e^{D_*}+2\sinh D_*=-e^{-D_*}.
\]
For \(d<0\) they already have the decaying signs. The remaining errors are exponentially small uniformly for \(|d|\) in a sufficiently small fixed interval, since the next height rate is \(4\) and \(3+2|d|<4\). The limiting quadratic form in a core vector \(z\) and added coordinate \(r\) is
\[
(z+e^{-D_*}e_F r)^tA_\sigma^\infty
(z+e^{-D_*}e_F r)+\tfrac12(1-e^{-2D_*})r^2,
\]
which has a strict gap when \(D_*\) is bounded below. Together the two regimes cover every frequency for sufficiently small fixed \(|d|>0\).

For the odd directions the outer part alone is
\[
\begin{pmatrix}
o(0)&o(1/2+d)\\
o(1/2+d)&(o(0)+o(1+2d))/2
\end{pmatrix}\ge c_*\min(k,1)I
\]
for small \(|d|\), by strict Green positivity and its limits at zero and infinity. The short part is nonnegative. The pre-collision \(H\)-part vanishes. In the post-collision case its uncontinued part is the positive rank-one form from the two sinh moments above. The only adverse continuation is the last entry of \eqref{j:height-continuations}. On bounded frequency intervals its cost is \(O(|d|k)\). At large \(k\), its sum with the leading height diagonal is
\[
e^{2D_*}-2\sinh(2D_*)=e^{-2D_*},
\]
up to exponentially small errors, and the height cross is exponentially small. First choose a large-frequency cutoff that absorbs these errors uniformly for \(|d|\) in a small fixed interval. Then choose \(|d|\) small enough on the remaining bounded frequency interval. This proves the odd gap.

\paragraph{The physical limit and \(J\).}
The sector gaps show that the only null directions are the two \(O\) relations. At small \(k\) the common channel supplies the \(\boldsymbol c\boldsymbol c^t/k\) term, and all quotient directions have a \(k\)-order gap. After physical congruence the same assertions hold with the physical count and charge. More precisely, the expansions of \(p,o,\mathsf h_\sigma\), the continued term, and \(K\) give
\[
\mathcal H'(k)=\frac{a_c}{k}\boldsymbol c\boldsymbol c^t+kB_d+O_d(k^2)
\]
for a real symmetric matrix \(B_d\). This records the exact leading coefficient that will cancel the state square. For fixed nonzero \(d\), the preceding high-frequency formulas give exponential convergence to a positive quotient matrix.

The physical high-frequency limit is also entrywise nonnegative. On the same-species \(D\) labels \(P,M,O_+,O_-\), the endpoint coincidences give
\[
\begin{pmatrix}
1&0&1/2&1/2\\
0&1&1/2&1/2\\
1/2&1/2&1&0\\
1/2&1/2&0&1
\end{pmatrix}.
\]
The common outer endpoint compensates each opposite short or outer coincidence. The \(U,L\) diagonals are \(1\), with no other surviving \(U,L\) entries for fixed nonzero \(d\). In normalized fused coordinates, the \(F,F\) diagonal is \(1/4\) and the cross-species \(F,F\) entry is \(1/2\); these come from the extreme height difference \(3\). There are no other surviving fused crosses. Physical congruence changes these last two values to \(1\) and \(2\), preserving nonnegativity. Restriction gives the same conclusion on the actual pre-collision labels.

Finally choose \(2e\) below the fixed exponential convergence rate of \(\mathcal H'\) at infinity. There the positive term
\(\mathcal H'(\infty)(1-\tanh(ek))\) dominates the convergence error, so \(J\) has a positive exponential lower bound on the quotient. At zero the subtracted term is \(O(ek)\); reduce \(e\) to retain the low gap. Reducing it again on a remaining compact frequency interval preserves positivity there as well. The resulting low, intermediate, and high bounds imply \eqref{j:stability-J}: at high frequency the positive exponential divided by \(k\) dominates \(e^{-k^2}\), and at low frequency the two required rows are controlled by the quotient gap. Both \(\mathfrak n\) and \(\boldsymbol c\) annihilate \(\mathcal N\), so the full-coordinate inequality follows.
\end{proof}

\subsection{The Fourier energy and box confinement}

We now prove Theorem~\ref{j:absolute-stability}. Split the full matrix as
\[
\mathcal H(k)
=Q_\infty\tanh(ek)
+\bigl(K(k)-K(\infty)\tanh(ek)\bigr)+J(k).
\]
After multiplication by \(2\pi/k\), each of the first two matrices represents an entrywise nonnegative pair potential. To see this, put
\[
\phi_L(x)=-2\log|\tanh(\pi x/(4L))|.
\]
It is positive and increasing in \(L\), and its transform is \(2\pi\tanh(Lk)/k\). The first matrix has nonnegative coefficients by the lemma. Each nonzero entry of \(K\) is a positive multiple of
\(\tanh(2k)e^{-rk}\), with \(r\ge0\); here \(r>0\) exactly when an entry involves \(U\) or \(L\). If \(r=0\), the difference above gives a positive multiple of \(\phi_2-\phi_e\). If \(r>0\), \(K(\infty)\) vanishes in that entry and the potential is the convolution of \(\phi_2\) with the positive Poisson kernel of parameter \(r\). These pair potentials can therefore be dropped from a lower bound on the negative log weight, on every forward pair and self-image.

For the remaining regular part, use the even symbol
\[
\widehat{\mathcal R}(k)
=\frac{2\pi}{|k|}
\left(J(|k|)-\frac{a_c}{|k|}\boldsymbol c\boldsymbol c^t\right),
\qquad k\ne0,
\]
with its continuous value at \(k=0\). The low expansions above show bounded one-sided derivatives there; terms in \(|k|\) are allowed. At infinity the symbol is integrable, with tail \(O(k^{-2})\), and its piecewise second derivatives are integrable. Fourier inversion and two integrations by parts therefore show that \(\mathcal R(x)\) is continuous, bounded, and \(O((1+x^2)^{-1})\). Its periodization
\[
\mathcal R_T(x)=\sum_{r\in\mathbb Z}\mathcal R(x+rT)
\]
converges absolutely.

Index the physical particles by \(\alpha=1,\ldots,N\), with labels \(\iota_\alpha\) and positions \(x_\alpha\) in one period, and define the vector of point measures
\[
\rho_\iota=\sum_{\alpha:\,\iota_\alpha=\iota}\delta_{x_\alpha},
\qquad
\widehat\rho(k)=\int_T e^{-ikx}\rho(dx),\qquad
k\in(2\pi/T)\mathbb Z.
\]
The forward regular energy, including nonzero self-images, is
\[
\begin{aligned}
E_{\mathcal R}
&:=\sum_{\alpha,\beta}\ 
\sum_{\substack{r\in\mathbb Z\\x_\beta+rT>x_\alpha}}
\mathcal R_{\iota_\alpha\iota_\beta}(x_\beta-x_\alpha+rT)\\
&=\frac12\sum_{\alpha,\beta}
(\mathcal R_T)_{\iota_\alpha\iota_\beta}(x_\beta-x_\alpha)
-\frac12\sum_\alpha\mathcal R_{\iota_\alpha\iota_\alpha}(0)\\
&=\frac1{2T}\sum_{k\in(2\pi/T)\mathbb Z}
\widehat\rho(k)^*\widehat{\mathcal R}(k)\widehat\rho(k)
-\frac12\sum_\alpha\mathcal R_{\iota_\alpha\iota_\alpha}(0).
\end{aligned}
\]
Indeed the half-double sum pairs the forward image \((\alpha,\beta,r)\) with its reverse \((\beta,\alpha,-r)\). Only \((\alpha,\alpha,0)\) is excluded from the gas. The subtracted diagonal is therefore the unperiodized value \(\mathcal R_{\iota_\alpha\iota_\alpha}(0)\); every nonzero self-image remains.

Closure of the state path gives
\[
\boldsymbol c^t\widehat\rho(0)=0,\qquad
p_c'=-\boldsymbol c^t\rho,\qquad
ik\widehat p_c(k)=-\boldsymbol c^t\widehat\rho(k)\quad(k\ne0).
\]
The common component \(p_c=(p_1+p_2)/\sqrt2\) is periodic because the total charge is in \(8\mathbb Z(1,-1)\). Put \(\eta=\theta/(\sqrt2\pi)\) and
\(\bar p_c=T^{-1}\int_Tp_c\). Parseval's identity now gives the exact cancellation
\begin{equation}\label{j:stability-energy}
\begin{aligned}
E_{\mathcal R}
+\pi a_c\Re\int_T(p_c+\eta)^2\,dx
={}&\frac{\pi}{T}\sum_{k\ne0}
\widehat\rho(k)^*\frac{J(|k|)}{|k|}\widehat\rho(k)
+\frac1{2T}\widehat\rho(0)^t\widehat{\mathcal R}(0)\widehat\rho(0)\\
&+\pi a_cT\bigl[(\bar p_c+\Re\eta)^2-(\Im\eta)^2\bigr]
-\frac12\sum_\alpha\mathcal R_{\iota_\alpha\iota_\alpha}(0).
\end{aligned}
\end{equation}
Here and below the sum uses \(k\in(2\pi/T)\mathbb Z\). At each nonzero frequency the negative term
\(-\pi a_c|\boldsymbol c^t\widehat\rho(k)|^2/(Tk^2)\) in the pair energy is canceled by the state square. At zero, \(\widehat\rho(0)\) is perpendicular to \(\boldsymbol c\), so the limit of \eqref{j:stability-J} gives
\[
\frac1{2T}\widehat\rho(0)^t\widehat{\mathcal R}(0)\widehat\rho(0)
\ge\frac{\pi\gamma}{T}
\left|\mathfrak n^t\widehat\rho(0)\right|^2.
\]
No positivity of \(\widehat{\mathcal R}(0)\) on other vectors is required. The mean square in \eqref{j:stability-energy} bounds a positive multiple of \(T\bar p_c^2\) up to \(O_\Theta(T)\). The diagonal correction costs \(O_d(N)\), independently of \(T\).

Let \(g_T\) be the normalized periodic Gaussian whose Fourier multiplier is \(e^{-k^2/2}\), and let
\(\nu=\mathfrak n^t\rho=\sum_\alpha\delta_{x_\alpha}\). Combining \eqref{j:stability-J} with \eqref{j:stability-energy}, and including the local factors, gives for every nonzero post-collision weight
\[
-\log|W_{d,\theta}|
\ge c_2\left(\|g_T*\nu\|_{L^2(T)}^2+\|g_T*p_c\|_{L^2(T)}^2\right)
-C_2(N+T).
\]
Here \(c_2,C_2\) depend only on \(d,\Theta\). The local constants have logarithmic cost \(O_d(N)\). The Berry factors containing real state charges are phases; their \(\theta\)-dependent parts cost \(O_\Theta(N)\). For the pre-collision product \eqref{j:pre-weight}, use the same identity with the real angle \(\vartheta=\Re\theta\). Its symmetrized local constants are finite for fixed \(d<0\), so they cost \(O_d(N)\), and the same bound holds with \(-\log W^{\rm pre}_{d,\Re\theta}\) on the left. The explicit heat and deterministic costs outside that product are absorbed afterward as described following the theorem. Zero weights already satisfy the desired inequality.

The first norm controls the box counts. In fact, with \(h_T=g_T*g_T\),
\[
\|g_T*\nu\|_2^2=\sum_{\alpha,\beta}h_T(x_\alpha-x_\beta)
\ge c_1\sum_{j=0}^{T-1}n_j^2,
\]
because \(h_T\) is positive and has a uniform positive lower bound for pairs in the same unit box. Each jump of \(p_c\) has magnitude at most \(\sqrt2\). For \(x\) in box \(j\), express \(p_c(x)-(g_T*p_c)(x)\) by integrating its increments against the Gaussian. A jump in a box \(r\) boxes away is then multiplied by a Gaussian tail bounded by \(C e^{-c(r-2)_+^2}\). The resulting sequence of bounds is summable uniformly after periodization. Discrete Young's inequality therefore gives
\[
\|p_c-g_T*p_c\|_2^2\le C_1\sum_j n_j^2.
\]
Within box \(j\), the same jump bound gives
\[
\sum_j p_{c,j}^2\le 2\|p_c\|_2^2+C_2\sum_j n_j^2.
\]
Using a smaller positive constant, the two smoothed norms consequently control
\(\sum_j(n_j^2+p_{c,j}^2)\). Finally
\(\sum_j n_j\le\varepsilon\sum_jn_j^2+C_\varepsilon T\)
absorbs the linear costs. This proves Theorem~\ref{j:absolute-stability}.

The pointwise bound gives the required absolute summability. For a fixed common integer \(p_1+p_2\), there are eight classes in \(\mathbb Z^2/8\mathbb Z(1,-1)\); the Gaussian bound on the common charge sums over that integer. In each box, integration and the ordinary particle factorials reduce to a convergent sum of the form \(\sum_{n\ge0}A^n e^{-cn^2}/n!\). Thus both the pre-collision Gaussian integration and the post-collision gas series are absolutely summable at the fixed displacements used in Section~\ref{c7:boson}.
